\section{Related Work}
\label{sec:related_work}

\paragraph{Temporal abstraction in RL.}
\looseness=-1 Temporal abstraction spans hierarchical, action-repetition, and world-model approaches, none of which has examined the interaction between flat $\dt$-conditioning and imagination-based value estimation.
The options framework \citep{sutton1999options}, option-critic \citep{bacon2017optioncritic}, HAM \citep{parr1998ham}, and MAXQ \citep{dietterich2000maxq} formalize temporal abstraction through hierarchical policies. Semi-MDP theory \citep{bradtke1995smdp} provides the mathematical foundation. Action-repetition methods achieve simpler temporal abstraction via real transitions: FiGAR \citep{sharma2017figar} learns per-action repetition counts jointly with the policy, TempoRL \citep{biedenkapp2021temporl} and \citet{park2021timediscretization} provide complementary formulations. SPiRL \citep{pertsch2020spirl} accelerates RL through learned skill priors that capture reusable temporal abstractions from offline data. \citet{wang2024elastic} learn elastic time steps that adapt step size to state dynamics. The Dreamer family \citep{hafner2020dreamerv1, hafner2021dreamerv2, hafner2025dreamerv3, hafner2025dreamerv4} trains policies in imagination at fixed $\dt{=}1$; hierarchical variants \citep{hafner2022director, gumbsch2024thick, mattes2024hieros, li2025lsimagine} use subgoals or separate dynamics per scale, avoiding flat $\dt$-conditioning. TAWM \citep{nhu2025tawm} conditions a world model on $\dt$ within TD-MPC2 \citep{hansen2024tdmpc2}; MPC re-plans from real states, avoiding the critic accumulation we analyze (\S\ref{sec:discussion}).

\paragraph{Value overestimation and model fidelity.}
\looseness=-1 Overestimation has been studied extensively in model-free and model-based RL, but no prior work identifies the specific pathology arising from learned temporal conditioning in imagination.
Maximization bias \citep{thrun1993issues} has been addressed by Double Q-learning \citep{vanhasselt2016double}, TD3 \citep{fujimoto2018td3}, and ensemble methods \citep{chen2021redq}. Scaling to larger critic ensembles is unlikely to help in pure imagination: our dual-critic experiments show $\rho > 0.999$ convergence regardless of capacity, because critics share the same imagined training distribution. In model-based settings, MBPO \citep{janner2019mbpo} and MOPO \citep{yu2020mopo} bound rollout error, while MOReL \citep{kidambi2020morel} and COMBO \citep{yu2021combo} penalize model uncertainty. Unlike MOPO, which penalizes model uncertainty using real-data rollouts as an anchor, pure imagination lacks this external correction mechanism, allowing bias to compound unchecked. \citet{lu2018nondelusional} identify delusional bias from policy-class constraints. CIVO is mechanistically distinct: the amplifier is the world model's co-adaptation with a learned temporal embedding, producing systematic reward-prediction bias at $\dt > 1$. Multi-timestep world models \citep{benechehab2023multitimestep} learn predictions at variable horizons but do not examine the value overestimation pathology we isolate. Compounding error in multi-step rollouts \citep{lambert2020objective, asadi2019combating} is a related but distinct phenomenon: if CIVO were generic compounding error, untrained $\dt$-embeddings at $\dt{=}2$ should also produce overestimation, yet they do not (NVE${<}0.6\times$ for unexposed and frozen controls vs.\ $4.12\times$ for the trained configuration). Signed reward error analysis reinforces this distinction: $\dt{=}2$ rewards are \emph{under}predicted (mean signed error [SE]${=}{-}0.599$), which rules out observation-level overprediction and localizes CIVO to imagination-level prior drift (\S\ref{sec:discussion}). Decision-aware model learning \citep{hamrick2021planning} optimizes for planning performance but does not address $\dt$-conditioned co-adaptation. Our temporal conditioning noise (TCN) intervention relates to training-with-noise \citep{bishop1995noise}; TCN's environment-specific dose sensitivity (\S\ref{sec:remedies}) suggests that CIVO reflects a structural vulnerability rather than a regularization deficiency. Our factorial diagnostic methodology follows \citet{henderson2018deep}.
This paper fills the gap at the intersection of temporal conditioning and imagination-based value estimation, revealing that shared imagined rollouts create a structural constraint: remedies that rely on training-data diversity (Double Q pessimism, ensembles, noise injection) lose their corrective power when both the training distribution and bootstrap targets are generated by the same learned model.
